\section{BabyLM：用数据约束反问规模崇拜}

第三个单元把视角从 inference 转回 pretraining。BabyLM Challenge 不追求最大模型，而是限制在近似儿童语言暴露量的数据预算内，比较哪些 architecture、objective 和 curriculum 能更高效学习。它既是 sample-efficiency benchmark，也是对研究资源集中和人类语言学习机制的反思。

\subsection{Challenge 的核心约束}

标题页标记从 reasoning methods 转向 data efficiency。课堂背景是大模型训练成本持续上升、个人与大学实验室难以参与 frontier pretraining，因此通过固定小数据预算重新打开可比较研究空间。

\lecturefigure{slide-24-babylm-challenge.jpg}{BabyLM Challenge 以小数据预算研究高效语言学习}{Stanford 课堂视频 00:15:15}

\begin{knowledgebox}{读图：约束本身改变研究问题}
普通 scaling benchmark 问“更多 data/params 能到多高”；BabyLM 问“在 data cap 下怎样用得更好”。这会奖励 inductive bias、objective、tokenization、curriculum 和 architecture，而不只奖励获取更多网页与算力。
\end{knowledgebox}

\subsection{什么是 developmentally plausible data budget}

课堂把 BabyLM 描述为：训练较小语言模型，使用接近儿童成长期间可接触的语言数据量。Slide 对比 Chinchilla 级别万亿词与儿童每年约数百万词的数量差距，目标是在显著更少数据下保持语法、语义和推理能力。本节先解释这一数量级约束，再说明它为何只是 developmentally inspired，而不是完整儿童学习模型。

\lecturefigure{slide-25-what-is-babylm.jpg}{BabyLM 对比大模型语料量与儿童语言暴露量}{Stanford 课堂视频 00:15:48}

\begin{knowledgebox}{读图：数据量对比是数量级启发}
儿童词汇暴露估计依赖年龄、家庭与统计方法；LLM token 也不等于人类听到的 word。这个对比用于强调现代预训练样本效率很低，不应用来声称文本模型和儿童拥有相同感知、互动或学习目标。
\end{knowledgebox}

\begin{importantbox}{样本效率应连同预算报告}
若模型在任务集上的平均性能为 $S$，训练 token 数为 $N$，可用
\[
\eta=\frac{S-S_0}{\log(1+N)}
\]
作粗略 efficiency 视角，其中 $S_0$ 是随机或小模型基线。$\eta$ 不是 BabyLM 官方指标，只说明比较时必须同时报告 performance 与 data budget，而不能只看绝对分数。
\end{importantbox}

\subsection{为什么值得做 BabyLM}

Why BabyLM slide 列出效率、开放新用例、interpretability、alignment、university access、认知科学和人类语言学习。它把小模型研究从“做不起大模型的替代品”转成独立科学问题：什么数据与结构能产生高效 generalization。

\lecturefigure{slide-26-why-babylm.jpg}{BabyLM 的效率、可及性与跨学科动机}{Stanford 课堂视频 00:17:00}

\begin{knowledgebox}{读图：七个理由分属工程与科学两类}
工程侧关注训练成本、部署与大学预算；科学侧关注可解释性、alignment、cognitive science 和 child language acquisition。小模型更容易做 controlled ablation，但参数少不自动意味着可解释或安全。
\end{knowledgebox}

\teachervoice{讲者强调，大模型研究越来越受限于拥有数百上千 GPU 和数百万美元的机构。BabyLM 的价值之一，是让大学和个人重新能够训练完整模型、比较假设，而不只调用闭源 API。}

\subsection{数据组成与儿童类比的边界}

最终 data slide 展示 developmentally plausible pretraining mixture，包括 CHILDES child-directed speech、OpenSubtitles、儿童书籍、BNC speech、Wikipedia 与 Simple Wikipedia 等来源。不同数据源提供对话、故事、日常语法和百科知识，但仍是离线文本集合。

\lecturefigure{slide-27-babylm-data.jpg}{BabyLM 的 developmentally inspired pretraining data mixture}{Stanford 课堂视频 00:19:09}

\begin{knowledgebox}{读图：mixture 在模拟输入类型而非完整发展环境}
Child-directed speech 提供互动语言风格，subtitles 提供对话，books 提供叙事，Wikipedia 提供事实与长句。真正儿童还拥有视觉、行动、共同注意、纠正与社会反馈；文本 mixture 只能约束语言输入量，不能复现 embodied learning。
\end{knowledgebox}

\begin{warningbox}{Developmentally plausible 不等于 cognitively equivalent}
模型接受全量语料的训练 objective、epoch 和顺序与儿童持续学习不同；它也缺少身体和社会目标。BabyLM 能检验 sample-efficient language modeling，但不能单独证明人类语言习得理论。
\end{warningbox}

\teachervoice{讲者希望认知科学与 NLP 双向交流：人类高效学习可以启发模型，受控小模型实验也可能反过来检验语言学习假设。这里的价值是可比较实验平台，而不是把儿童简化为 next-token predictor。}

\subsection{本章小结}

BabyLM 用 sub-100M-word 级别预算研究 sample efficiency，推动 architecture、objective、curriculum 与数据 composition 比较，并降低研究门槛。儿童类比提供数量与类型启发，但不包含完整感知、互动和发展机制。

